\section{Video Experimental Setup Details}\label{sec:video_experimental_setup}

\subsection{Dataset Details}\label{subsec:video_dataset}

We train on 178,688 and evaluate on 2,048 clips (disjoint) of 32 frames each from CelebV-HQ \cite{zhu2022celebv}. We train on 17,152 and evaluate on 1,360 clips (disjoint) of 64 frames each from Sky-Timelapse \cite{zhang2020dtvnet}. We diffuse in the Stable Diffusion latent space \cite{rombach2022high} on 32x32x4 codes, corresponding to 256x256x3 images.

\subsection{Training}\label{subsec:training_details}

We sample the time $t \sim \mathcal{U}[0, 1]$, then clip it so that $|t - t_i| > 10^{-4}$, for every finite-dimensional marginal time $t_i$. For the autoregressively conditioned diffusion bridge (and noise-to-data diffusion), we get the auxiliary time $\tau_i = \frac{t - t_{i^*}}{t_{i^*+1} - t_{i^*}}$, as constructed in Equation \ref{eqn:time_shifted_bridge}.

We have cosine decay learning rate, with linear warmup for 10K steps up to rate $10^{-4}$, then decay for 290K steps. Due to compute limits, we evaluate the model from 240K steps. We use Adam \cite{kingma2014adam} without weight decay, and $\beta = (0.9, 0.999)$. We clip the gradient to maximum norm of $5.0$. We use batch size $256$. We run each training run on 4 A100 80GB GPUs, and each run takes about 36 hours.

\subsection{SDE Parameter Selection}\label{subsec:sde_parameter_selection}

There is a very large design space of SDE parameters that could have been used in Equations \ref{eqn:new_sde} and \ref{eqn:cond_diffusion_bridge}. To isolate the effects of \rqref{time_scaling} (impact of time-dependent quadratic variation), \rqref{data_to_data} (impact of data-to-data bridge), we picked relatively simple choices. Future work can explore more complicated parameterizations.

We always set $t_L = 1$, ensuring that the SDE is defined on $[0, 1]$.

We set the base (non-score driven) drift to be $a(t) = 0$. The score network should be able to learn the drift, without need for an explicit specification.

We use constant volatility $\sigma(t) = \sigma$ in ABC and the conditional diffusion bridge baseline.
This is because a non-constant volatility (say, small at the finite-dimensional marginal times, but large in between them, as is popular in DDBM \cite{zhou2023denoising}) would mean that the generative process would not have physically meaningful intermediate values in between the times at which we observe the finite-dimensional marginals. This is a non-starter for continuous-time modeling, as it kills any hope of generalization to modeling times that were not on the grid that the training data was sampled from.
In contrast, using a constant volatility decouples the sampling rate of the data from the underlying model of the stochastic process.

We selected the volatility by testing reconstruction fidelity in the Stable Diffusion VAE latent space \cite{rombach2022high} with various Gaussian noise perturbations. We found that perturbations up to $\mathcal{N}(0, 0.1^2 I)$ minimally affected the reconstruction quality. Assuming simulation step sizes up to $\Delta t \leq 0.01$, the upper limit on the volatility introduced from Brownian motion should be $\sigma \sqrt{\Delta t} \leq 0.1 \Rightarrow \sigma \leq 1.0$.

To select volatility for the baseline conditional diffusion bridge, we tried a few options.
The first, most naive one was to set the volatility coefficient exactly the same as that of ABC's:
\begin{equation}
    \sigma_\text{cond bridge} = \sigma_\text{ABC}
\end{equation}
Another option was to calculate the volatility coefficient based on the quadratic variation (up to rounding) that would have been introduced by ABC in between frames (given a frame sampling rate):
\begin{equation}
    \sigma_\text{cond bridge} = \sqrt{t_{i^*+1} - t_{i^*}} * \sigma_\text{ABC},
\end{equation}
in accordance with time-scaling rules for Brownian motion \cite{steele2001stochastic}.
Finally, when compute resources permitted, we tried a few other volatility values.

For the noise-to-data diffusion baseline, we used exponentially decaying volatility (Section \ref{subsec:exponential_decay_volatility}) and cosine decaying volatility (Section \ref{subsec:cosine_decay_volatility}).

\subsection{Inference}

We also concatenate the final frame from $t = 1$ as a prompt.

\subsection{Evaluation}\label{subsec:evaluation_details}

FVD and FID are computed on the infilled videos (including the prompt frames), with the same number of videos (1360 in Sky Timelapse, 2048 in CelebV-HQ) for each of the ground truth and fake sets. We exclude the final frame ($t=1$) from computations.
FVD computations use VideoMAE v2 backbone with interpolated time embeddings \cite{wang2023videomae, tong2022videomae, ge2024content}. FID computations use the default Inception net \cite{szegedy2016rethinking}.
